%% toc: P3 39 39.7
\chapter{Datalog, Existential Pebble Games and Fixed-Point Logics in Finite-Variable Infinitary Logic}
\label{sec:p03-w03-p03-060-datalog-infinitary}

This chapter places three recursive formalisms on finite structures inside
the finite-variable infinitary logics. Its inputs are a vocabulary, the
\(k\)-variable infinitary language \(\mathrm{L}^k_{\infty\omega}\) with
variable reuse (\cref{def:p03-w03-p03-056-grammar}), the fixed-point logics
LFP and PFP (\cref{def:p03-w03-p03-059-lfp,def:p03-w03-p03-059-pfp}) and
Datalog (\cref{def:p09-datalog}). Its conclusions are:
\begin{enumerate}
\item an existential \(k\)-pebble game characterizes preservation of
      existential \(k\)-variable infinitary formulas, in a version without
      and a version with inequalities
      (\cref{thm:p03-w03-p03-060-existential-game});
\item every LFP, IFP and PFP formula is equivalent on finite structures to a
      formula of \(\mathrm{L}^\omega_{\infty\omega}\)
      (\cref{thm:p03-w03-p03-060-fp-in-lomega}); for a single fixed point
      of a first-order formula with \(k\) variables, every finite stage is
      an \(\mathrm{FO}^k\)-formula on all structures and the fixed point an
      \(\mathrm{L}^k_{\infty\omega}\)-formula on finite structures, with the
      same \(k\) (\cref{thm:p03-w03-p03-060-kv-stages});
\item every Datalog query is equivalent on all structures to a formula of
      the existential fragment, and every Datalog query with inequalities to
      a formula of the existential fragment with inequalities
      (\cref{thm:p03-w03-p03-060-datalog-in-existential}); with the same
      number \(k\) of variables, the stages of a system of positive
      \(\exists\mathrm{FO}^k\)-formulas are \(\exists\mathrm{FO}^k\)-formulas
      and \(k\)-Datalog\((\neq)\) is contained in
      \(\exists\mathrm{L}^k_{\infty\omega}(\neq)\)
      (\cref{thm:p03-w03-p03-060-kv-systems,thm:p03-w03-p03-060-kdatalog});
\item every Datalog query is equivalent on finite structures to an LFP
      formula using only simple fixed points
      (\cref{thm:p03-w03-p03-060-datalog-in-lfp}); Datalog programs and
      least fixed points of existential-conjunctive systems define the same
      queries, with and without inequalities
      (\cref{thm:p03-w03-p03-060-datalog-systems}), and Datalog is equivalent
      to existential fixed-point logic on all structures
      (\cref{thm:p03-w03-p03-060-efp-datalog});
\item least, inflationary and partial fixed points are ordered
      \(\mathrm{LFP}\le\mathrm{IFP}\le\mathrm{PFP}\) on finite structures
      (\cref{thm:p03-w03-p03-060-fp-comparison}); inflationary stage
      comparison is LFP-definable, even in positive LFP, on finite
      structures (\cref{thm:p03-w03-p03-060-stage-comparison,thm:p03-w03-p03-060-stage-comparison-positive}),
      so IFP collapses to
      LFP there (\cref{thm:p03-w03-p03-060-gurevich-shelah}); and
      equivalence of LFP or IFP with PFP on all finite structures holds
      if and only if \(\mathsf P=\mathsf{PSPACE}\)
      (\cref{thm:p03-w03-p03-059-abiteboul-vianu}).
\end{enumerate}
The \(k\)-pebble game for full \(\mathrm{L}^k_{\infty\omega}\) and its
equivalence theorem are \cref{def:p03-w03-p03-056-pebble-game} and
\cref{thm:p03-w03-p03-056-eternal}; the existential games below are their
one-sided counterparts. Where a statement depends on the version of
Kolaitis's survey, it follows the manuscript of 1 August 2003 (a later
version is a chapter of the 2007 Springer volume); Gr\"adel's chapter is
cited from a copy in book layout posted on his institute web page, whose identity with the
printed text has not been checked.

\section{Existential finite-variable languages}
\label{sec:p03-w03-p03-060-languages}

\begin{hbconvention}{def:p03-w03-p03-060-setting}{Setting and homomorphisms}
Vocabularies \(\sigma\), \(\sigma\)-structures \(\mathbf A\) with nonempty
universe \(A\), the variables \(v_1,\ldots,v_k\), free labels
\(\operatorname{fr}(\varphi)\), assignments \(\alpha\colon D\to A\) with
\(D\subseteq[k]\), and satisfaction \(\mathbf A\models\varphi[\alpha]\) are
those of
\cref{def:p03-w03-p03-056-vocabulary,def:p03-w03-p03-056-grammar,def:p03-w03-p03-056-semantics};
\(\sigma\) is finite and may contain constant symbols. A \emph{partial
homomorphism} from \(\mathbf A\) to \(\mathbf B\) is a map \(h\) with
\(\{c^{\mathbf A}:c\in\sigma\}\subseteq\dom h\subseteq A\),
\(h(c^{\mathbf A})=c^{\mathbf B}\) for every constant symbol \(c\), and
\(R^{\mathbf A}(a_1,\ldots,a_r)\Rightarrow R^{\mathbf B}(h(a_1),\ldots,h(a_r))\)
for every \(r\)-ary \(R\in\sigma\) and \(a_1,\ldots,a_r\in\dom h\). It is
\emph{one-to-one} when \(h\) is injective. A \emph{homomorphism} is a
partial homomorphism with domain \(A\) \cite[Def.~6.22,
p.~44]{Kolaitis2003ExpressivePowerNotes}.
\end{hbconvention}

\begin{hbdefinition}{def:p03-w03-p03-060-existential}{The logics \(\exists\mathrm{L}^k_{\infty\omega}\), \(\exists\mathrm{L}^k_{\infty\omega}(\neq)\) and their first-order parts}
Let \(k\ge1\).
\begin{enumerate}
\item \(\exists\mathrm{L}^k_{\infty\omega}(\sigma)\) is the set of formulas
      generated from the atomic formulas of clause (G1) of
      \cref{def:p03-w03-p03-056-grammar} (equalities \(t=t'\) included) by
      \(\bigwedge\Phi\) and \(\bigvee\Phi\) for arbitrary sets \(\Phi\) and
      by \(\exists v_i\) for \(i\in[k]\); negation and \(\forall\) are not
      used.
\item \(\exists\mathrm{L}^k_{\infty\omega}(\neq)(\sigma)\) is generated in
      the same way from the atomic formulas together with the negated
      equalities \(t\neq t'\) between terms.
\item \(\exists\mathrm{FO}^k\) and \(\exists\mathrm{FO}^k(\neq)\) are the
      formulas of these sets in whose formation every \(\Phi\) is finite.
\item \(\exists\mathrm{L}^\omega_{\infty\omega}=\bigcup_{k\ge1}\exists\mathrm{L}^k_{\infty\omega}\)
      and
      \(\exists\mathrm{L}^\omega_{\infty\omega}(\neq)=\bigcup_{k\ge1}\exists\mathrm{L}^k_{\infty\omega}(\neq)\).
\end{enumerate}
For assignments \(\alpha\colon D\to A\), \(\beta\colon D\to B\) with a
common domain \(D\subseteq[k]\), write
\((\mathbf A,\alpha)\preceq^{\exists,k}_{\infty\omega}(\mathbf B,\beta)\)
when \(\mathbf A\models\varphi[\alpha]\) implies
\(\mathbf B\models\varphi[\beta]\) for every
\(\varphi\in\exists\mathrm{L}^k_{\infty\omega}\) with
\(\operatorname{fr}(\varphi)\subseteq D\);
\(\preceq^{\exists,k}_{\omega\omega}\) is the same with
\(\exists\mathrm{FO}^k\), and \(\preceq^{\exists,\neq,k}_{\infty\omega}\),
\(\preceq^{\exists,\neq,k}_{\omega\omega}\) use the \((\neq)\) languages. For
\(D=\emptyset\) the assignments are omitted.

These are Kolaitis's languages and relations \cite[Def.~8.1, p.~69 and
Def.~8.6, pp.~72--73]{Kolaitis2003ExpressivePowerNotes}, here over the
fixed variable supply of \cref{def:p03-w03-p03-056-grammar}. The empty
disjunction \(\bot\defeq\bigvee\emptyset\) and the empty conjunction
\(\top\defeq\bigwedge\emptyset\) belong to every one of these languages.
\index{existential infinitary logic}\index{finite-variable logic!existential}
\end{hbdefinition}

\section{Existential pebble games}
\label{sec:p03-w03-p03-060-games}

\begin{hbdefinition}{def:p03-w03-p03-060-existential-game}{The \((\exists,k)\)- and \((\exists,\neq,k)\)-pebble games}
Positions \((\alpha,\beta)\) with domain \(D\subseteq[k]\) and the
correspondence
\(\rho_{\alpha,\beta}=\{(\alpha(i),\beta(i)):i\in D\}\cup\{(c^{\mathbf A},c^{\mathbf B}):c\in\sigma\}\)
are as in \cref{def:p03-w03-p03-056-pebble-game}. In the
\emph{\((\exists,k)\)-pebble game} on \(\mathbf A\) and \(\mathbf B\) every
move of the Spoiler is made in \(\mathbf A\):
\begin{enumerate}
\item[(P\(_\exists\))] he chooses \(i\in[k]\) and \(a\in A\); the
      Duplicator chooses \(b\in B\); the new position is
      \((\alpha[i\mapsto a],\beta[i\mapsto b])\) (if \(i\in D\) the pebble
      pair \(i\) is lifted first);
\item[(R)] he chooses \(i\in D\); the new position is
      \((\alpha\restriction(D\setminus\{i\}),\beta\restriction(D\setminus\{i\}))\).
\end{enumerate}
A position is \emph{\(\exists\)-good} when \(\rho_{\alpha,\beta}\) is a
function and a partial homomorphism from \(\mathbf A\) to \(\mathbf B\);
it is \emph{\(\exists{\neq}\)-good} when \(\rho_{\alpha,\beta}\) is moreover
one-to-one. The Duplicator wins a play of the \((\exists,k)\)-pebble game
(of the \((\exists,\neq,k)\)-pebble game) when every position of the play,
the initial one included, is \(\exists\)-good (\(\exists{\neq}\)-good); a
play lasts forever, and strategies are as in
\cref{def:p03-w03-p03-056-pebble-game}.

An \emph{\(\exists\)-family} (\emph{\(\exists{\neq}\)-family}) is a set
\(\mathcal W\) of positions over \([k]\) such that every member is
\(\exists\)-good (\(\exists{\neq}\)-good), \(\mathcal W\) is closed under
the removal (R), and for every \((\alpha,\beta)\in\mathcal W\), \(i\in[k]\)
and \(a\in A\) there is \(b\in B\) with
\((\alpha[i\mapsto a],\beta[i\mapsto b])\in\mathcal W\).
\end{hbdefinition}

Kolaitis's game places only pebbles not on the board and removes pebbles by
separate moves; the payoff asks that the pebbled correspondence be a
homomorphism, or a one-to-one homomorphism, between the generated
substructures \cite[Def.~8.4, pp.~71--72]{Kolaitis2003ExpressivePowerNotes}.
His winning strategies are families of partial homomorphisms closed under
subfunctions with the forth property up to \(k\) \cite[Def.~8.5,
p.~72]{Kolaitis2003ExpressivePowerNotes}. Reusing a pebble is a removal
followed by a placement, so a strategy for his game yields one for the game
above and conversely; the theorem below is stated for the game above.

\begin{hbtheorem}{thm:p03-w03-p03-060-existential-game}{proved=annex:p03-w03-p03-060-existential-game}{Existential pebble games characterize existential \(k\)-variable preservation}
Let \(k\ge1\), let \(\mathbf A,\mathbf B\) be \(\sigma\)-structures and
\((\alpha,\beta)\) a position with domain \(D\subseteq[k]\). The following
are equivalent:
\begin{enumerate}
\item the Duplicator has a winning strategy in the \((\exists,k)\)-pebble
      game on \(\mathbf A\) and \(\mathbf B\) from \((\alpha,\beta)\);
\item some \(\exists\)-family contains \((\alpha,\beta)\);
\item \((\mathbf A,\alpha)\preceq^{\exists,k}_{\infty\omega}(\mathbf B,\beta)\).
\end{enumerate}
They imply \((\mathbf A,\alpha)\preceq^{\exists,k}_{\omega\omega}(\mathbf B,\beta)\),
and if \(\mathbf B\) is finite this is equivalent to items~1--3. The same
holds with the \((\exists,\neq,k)\)-pebble game,
\(\exists{\neq}\)-families, \(\preceq^{\exists,\neq,k}_{\infty\omega}\) and
\(\preceq^{\exists,\neq,k}_{\omega\omega}\).
\end{hbtheorem}

For \(D=\emptyset\) this is Kolaitis's Theorem~8.7, which he credits to
Kolaitis and Vardi and states without proof
\cite[Thm.~8.7, p.~73]{Kolaitis2003ExpressivePowerNotes}. The
finiteness hypothesis falls on \(\mathbf B\) alone: it makes the
conjunction over the elements of \(B\) in the extension step finite.

\begin{hbcorollary}{thm:p03-w03-p03-060-existential-definability}{proved=annex:p03-w03-p03-060-existential-definability}{Existential definability on classes of finite structures}
Let \(\mathcal C\) be a class of finite \(\sigma\)-structures and
\(\mathcal K\subseteq\mathcal C\) closed under isomorphism. Then
\(\mathcal K\) is definable on \(\mathcal C\) by a sentence of
\(\exists\mathrm{L}^\omega_{\infty\omega}\) if and only if there is
\(k\ge1\) such that for all \(\mathbf A,\mathbf B\in\mathcal C\): if
\(\mathbf A\in\mathcal K\) and the Duplicator wins the
\((\exists,k)\)-pebble game on \(\mathbf A\) and \(\mathbf B\), then
\(\mathbf B\in\mathcal K\). The same holds for
\(\exists\mathrm{L}^\omega_{\infty\omega}(\neq)\) with the
\((\exists,\neq,k)\)-pebble games.
\end{hbcorollary}

This is Kolaitis's Corollary~8.8, stated there for
\(\exists\mathrm{L}^\omega_{\infty\omega}\) \cite[Cor.~8.8 and Method~8.9,
p.~73]{Kolaitis2003ExpressivePowerNotes}.

\begin{hbcorollary}{thm:p03-w03-p03-060-preservation}{proved=annex:p03-w03-p03-060-preservation}{Preservation under homomorphisms}
Let \(h\colon\mathbf A\to\mathbf B\) be a homomorphism and
\(\varphi\in\exists\mathrm{L}^\omega_{\infty\omega}\). For every assignment
\(\alpha\) in \(\mathbf A\) admissible for \(\varphi\),
\(\mathbf A\models\varphi[\alpha]\) implies
\(\mathbf B\models\varphi[h\circ\alpha]\). If \(h\) is one-to-one, the same
holds for every \(\varphi\in\exists\mathrm{L}^\omega_{\infty\omega}(\neq)\).
\end{hbcorollary}

Kolaitis records these preservation properties for the two logics
\cite[p.~71]{Kolaitis2003ExpressivePowerNotes} and the corresponding
ones for Datalog and Datalog\((\neq)\) as Propositions~6.23 and~6.26
\cite[pp.~44--45]{Kolaitis2003ExpressivePowerNotes}.

\begin{hbtheorem}{thm:p03-w03-p03-060-cliques}{proved=annex:p03-w03-p03-060-cliques}{The existential game is not symmetric}
Let \(\mathbf K_m\) be the complete graph on \(m\) vertices (no loops) and
\(k\ge1\). The Spoiler wins the \((\exists,k+1)\)-pebble game on
\(\mathbf K_{k+1}\) and \(\mathbf K_k\), and the Duplicator wins the
\((\exists,\neq,j)\)-pebble game on \(\mathbf K_k\) and \(\mathbf K_{k+1}\)
for every \(j\ge1\).
\end{hbtheorem}

The two halves are Kolaitis's illustration of the asymmetry
\cite[p.~72]{Kolaitis2003ExpressivePowerNotes}. With
\cref{thm:p03-w03-p03-060-existential-game}, the first half gives an
\(\exists\mathrm{FO}^{k+1}\)-sentence true in \(\mathbf K_{k+1}\) and false
in \(\mathbf K_k\), while by the second half no
\(\exists\mathrm{L}^\omega_{\infty\omega}(\neq)\)-sentence true in
\(\mathbf K_{k+1}\) and false in \(\mathbf K_k\) exists in the opposite
direction.

\section{Fixed-point logics inside \(\mathrm{L}^\omega_{\infty\omega}\)}
\label{sec:p03-w03-p03-060-fp}

\begin{hbdefinition}{def:p03-w03-p03-060-ifp}{Inflationary fixed points and IFP}
Let \(\mathbf A\) be a structure and \(\psi(R,\bar x)\) a formula with
\(R\) of arity \(r\), defining (with its parameters fixed) the operator
\(F_\psi(S)=\{\bar a:(\mathbf A,S)\models\psi(S,\bar a)\}\) on
\(\mathcal P(A^r)\). Its \emph{inflationary stages} are
\(R^0\defeq\varnothing\) and \(R^{m+1}\defeq R^m\cup F_\psi(R^m)\), and on a
finite structure the \emph{inflationary fixed point} is
\(\bigcup_{m\in\N}R^m\). IFP extends first-order logic by the rule: from any
formula \(\psi(R,\bar x)\), a tuple \(\bar x\) of distinct variables and a
tuple \(\bar t\) of terms of the arity of \(R\), form
\([\mathbf{ifp}\,R\bar x\,.\,\psi](\bar t)\), true when the value of
\(\bar t\) lies in the inflationary fixed point. No positivity is required
\cite[Def.~3.3.45, p.~174]{Gradel2007FMTDescriptiveChapter}. Free variables
and parameters are as in \cref{def:p03-w03-p03-059-lfp}.
\end{hbdefinition}

On a finite structure the inflationary stages increase and stop after at
most \(|A|^r\) proper increases, so the union is reached at a finite stage
and is a fixed point of \(S\mapsto S\cup F_\psi(S)\).

\begin{hblemma}{thm:p03-w03-p03-060-substitution}{proved=annex:p03-w03-p03-060-substitution}{Substituting a defined relation for a relation symbol}
Let \(R_1,\ldots,R_n\) be relation symbols of arities \(r_1,\ldots,r_n\), let
\(w\ge\max_ir_i\), and let \(x_1,\ldots,x_w,z_1,\ldots,z_w\) be distinct
variables. Let \(\chi\) be a formula of \(\mathrm{L}_{\infty\omega}\) with
finitely many variables, of vocabulary \(\tau\cup\{R_1,\ldots,R_n\}\), and
for each \(l\) let \(\theta_l\) be a formula of \(\mathrm{L}_{\infty\omega}\),
LFP, IFP or PFP of vocabulary \(\tau\) whose free variables outside
\(x_1,\ldots,x_{r_l}\) (its \emph{parameters}) are bound nowhere in
\(\chi\). Assume no \(z_j\) occurs in \(\chi\) and no \(z_j\) is free in
any \(\theta_l\).
Let \(\chi[\bar R:=\bar\theta]\) be obtained from \(\chi\) by replacing each
atomic subformula \(R_l(t_1,\ldots,t_{r_l})\) by
\[
 \exists z_1\cdots\exists z_{r_l}\Bigl(\textstyle\bigwedge_{j\le r_l}z_j=t_j\wedge
 \exists x_1\cdots\exists x_{r_l}\bigl(\bigwedge_{j\le r_l}x_j=z_j\wedge\theta_l\bigr)\Bigr).
\]
Then for every \(\tau\)-structure \(\mathbf A\) (finite when some
\(\theta_l\) uses \(\mathbf{pfp}\) or \(\mathbf{ifp}\)) and every assignment
\(s\) to the free variables of \(\chi\) and the parameters,
\(\mathbf A\models\chi[\bar R:=\bar\theta][s]\) if and only if
\((\mathbf A,\theta_1^{\mathbf A,s},\ldots,\theta_n^{\mathbf A,s})\models\chi[s]\),
where \(\theta_l^{\mathbf A,s}=\{\bar a\in A^{r_l}:\mathbf A\models\theta_l[s[\bar x\mapsto\bar a]]\}\).
The variables of \(\chi[\bar R:=\bar\theta]\) are among those of \(\chi\),
of the \(\theta_l\), and \(x_1,\ldots,x_w,z_1,\ldots,z_w\). If \(\chi\) and
every \(\theta_l\) lie in \(\exists\mathrm{L}_{\infty\omega}\) (with
\(\neq\)), so does \(\chi[\bar R:=\bar\theta]\); if \(R_l\) and a relation
symbol \(P\) occur only positively in \(\chi\) and \(P\) only positively in
\(\theta_l\), then \(P\) occurs only positively in
\(\chi[\bar R:=\bar\theta]\).
\end{hblemma}

The extra variables \(z_1,\ldots,z_w\) carry the argument values across the
requantification of \(x_1,\ldots,x_w\); without them an atom such as
\(R(x_2,x_1)\) could not be replaced by a copy of \(\theta\) written in
\(x_1,x_2\) without enlarging the variable set at every stage.

\begin{hbtheorem}{thm:p03-w03-p03-060-fp-in-lomega}{proved=annex:p03-w03-p03-060-fp-in-lomega}{LFP, IFP and PFP are contained in \(\mathrm{L}^\omega_{\infty\omega}\) on finite structures}
Let \(\varphi\) be a formula of LFP (with \(\mathbf{lfp}\) and
\(\mathbf{gfp}\), \cref{def:p03-w03-p03-059-lfp}), of IFP
(\cref{def:p03-w03-p03-060-ifp}) or of PFP (\cref{def:p03-w03-p03-059-pfp})
of vocabulary \(\sigma\), with \(P_1,\ldots,P_q\) its free relation
variables. There are \(k\) and a formula \(\varphi^\ast\) of
\(\mathrm{L}_{\infty\omega}(\sigma\cup\{P_1,\ldots,P_q\})\) with at most
\(k\) distinct variables, whose free variables are among those of
\(\varphi\), such that for every finite
\((\sigma\cup\{P_1,\ldots,P_q\})\)-structure and every assignment to the
free variables of \(\varphi\), \(\varphi\) and \(\varphi^\ast\) have the
same truth value. In particular, on every class \(\mathcal C\) of finite
\(\sigma\)-structures each query defined by an LFP, IFP or PFP formula is
defined by an \(\mathrm{L}^\omega_{\infty\omega}\)-formula.
\end{hbtheorem}

By \cref{thm:p03-w03-p03-056-renaming} a formula with at most \(k\) distinct
variables can be rewritten over \(v_1,\ldots,v_k\), so \(\varphi^\ast\) is,
up to that renaming, an \(\mathrm{L}^k_{\infty\omega}\)-formula. The
construction for nested formulas counts variables generously: if
\(\psi^\ast\) has \(k_\psi\) variables, \(\varphi^\ast\) has at most
\(k_\psi+|\bar p|+2r+|\operatorname{var}(\bar t)|\). For a single fixed
point of a first-order formula the count is exact: the stages use no
variable beyond those of the defining formula. The tool is a substitution
property of the \(k\)-variable languages.

\begin{hblemma}{thm:p03-w03-p03-060-variable-substitution}{proved=annex:p03-w03-p03-060-variable-substitution}{Substituting terms for variables in \(k\)-variable formulas}
Let \(k\ge1\) and let \(\mathcal L\) be one of
\(\mathrm{L}^k_{\infty\omega}\), \(\mathrm{FO}^k_{\omega\omega}\)
(\cref{def:p03-w03-p03-056-grammar}), \(\exists\mathrm{L}^k_{\infty\omega}\),
\(\exists\mathrm{FO}^k\), \(\exists\mathrm{L}^k_{\infty\omega}(\neq)\),
\(\exists\mathrm{FO}^k(\neq)\) (\cref{def:p03-w03-p03-060-existential}).
Let \(T_k\) be the set of terms over \(v_1,\ldots,v_k\) (the variables and
the constant symbols of \(\sigma\)). For every \(\varphi\in\mathcal L(\sigma)\)
and every map \(\rho\colon[k]\to T_k\) there is
\(\varphi_\rho\in\mathcal L(\sigma)\) such that
\[
 \operatorname{fr}(\varphi_\rho)\subseteq\{j:v_j\text{ occurs in }\rho(i)\text{ for some }i\in\operatorname{fr}(\varphi)\}
\]
and, for every \(\sigma\)-structure \(\mathbf A\) and every assignment
\(\alpha\colon[k]\to A\),
\[
 \mathbf A\models\varphi_\rho[\alpha]\iff\mathbf A\models\varphi[\rho^\alpha],
 \qquad\text{where }\rho^\alpha(i)\defeq\rho(i)^{\mathbf A}[\alpha]\ (i\in[k]).
\]
\end{hblemma}

For maps \(\rho\) into the variables this is the closure property that
Kolaitis isolates as the step to be proved separately for
\(\exists\mathrm{FO}^k\) \cite[proof of Thm.~8.2, p.~70]{Kolaitis2003ExpressivePowerNotes}.
The point is that \(\rho\) need not be injective: in
\(\varphi=\exists v_2\,E(v_1,v_2)\) with \(\rho(1)=v_2\) and \(\rho(2)=v_1\),
the bound variable must move out of the way of the substituted \(v_2\), and
the construction in the annex returns
\(\varphi_\rho=\exists v_1\,E(v_2,v_1)\), which says of the value of \(v_2\)
what \(\varphi\) says of the value of \(v_1\).

\begin{hbtheorem}{thm:p03-w03-p03-060-kv-stages}{proved=annex:p03-w03-p03-060-kv-stages}{Kolaitis and Vardi: the stages of a \(k\)-variable formula are \(k\)-variable formulas}
Let \(\sigma\) be a vocabulary, \(S\) an \(m\)-ary relation symbol not in
\(\sigma\), and \(\varphi(x_1,\ldots,x_m,S)\) a first-order formula of
vocabulary \(\sigma\cup\{S\}\) whose free variables are among the distinct
variables \(x_1,\ldots,x_m\). Fix a declared supply of \(k\ge1\) distinct
first-order variables containing both every variable occurring in
\(\varphi\), free or bound, and every coordinate variable
\(x_1,\ldots,x_m\); thus \(0\le m\le k\), and unused coordinates are
counted. On a \(\sigma\)-structure \(\mathbf A\) let
\(\Phi(P)=\{\bar a\in A^m:(\mathbf A,P)\models\varphi[\bar x\mapsto\bar a]\}\),
with stages \(\Phi^0=\varnothing\), \(\Phi^{n+1}=\Phi(\Phi^n)\),
inflationary stages \(\Phi^0_{\mathrm{inf}}=\varnothing\),
\(\Phi^{n+1}_{\mathrm{inf}}=\Phi^n_{\mathrm{inf}}\cup\Phi(\Phi^n_{\mathrm{inf}})\),
and downward stages \(\Phi_0=A^m\), \(\Phi_{n+1}=\Phi(\Phi_n)\). Rename the
declared supply bijectively to \(v_1,\ldots,v_k\) with
\(x_j\mapsto v_j\) (\cref{thm:p03-w03-p03-056-renaming}).
\begin{enumerate}
\item For every \(n\in\N\) there are formulas \(\varphi^n\),
      \(\varphi^n_{\mathrm{inf}}\), \(\varphi_n\) of
      \(\mathrm{FO}^k_{\omega\omega}(\sigma)\) with free labels in \([m]\)
      that define \(\Phi^n\), \(\Phi^n_{\mathrm{inf}}\), \(\Phi_n\) on every
      \(\sigma\)-structure, finite or infinite: for every
      \(\alpha\colon[k]\to A\),
      \(\mathbf A\models\varphi^n[\alpha]\iff(\alpha(1),\ldots,\alpha(m))\in\Phi^n\),
      and likewise for the other two.
\item On every finite \(\sigma\)-structure the partial fixed point of
      \(\varphi\) is defined by the \(\mathrm{L}^k_{\infty\omega}\)-formula
      \[
       \textstyle\bigvee_{n\in\N}\Bigl(\forall v_1\cdots\forall v_m\bigl((\neg\varphi^n\vee\varphi^{n+1})\wedge(\neg\varphi^{n+1}\vee\varphi^n)\bigr)\wedge\varphi^n\Bigr),
      \]
      the inflationary fixed point by
      \(\bigvee_{n\in\N}\varphi^n_{\mathrm{inf}}\), and, if \(S\) occurs
      only positively in \(\varphi\), the least fixed point by
      \(\bigvee_{n\in\N}\varphi^n\) and the greatest fixed point by
      \(\bigwedge_{n\in\N}\varphi_n\).
\end{enumerate}
\end{hbtheorem}

This is Kolaitis's Theorem~7.2, whose two items are the finite stages (for
\(n\ge1\); his first stage is \(\Phi(\varnothing)=\Phi^1\)) and the partial
fixed point; he refers to Kolaitis and Vardi for a detailed proof
\cite[Thm.~7.2, p.~54]{Kolaitis2003ExpressivePowerNotes}. The
inflationary and greatest fixed points are added here. The declared supply
also makes the coordinate convention explicit: when every coordinate
occurs in \(\varphi\), \(k\) may be its actual number of variables as in
the source; otherwise the unused coordinates must also fit in the supply.
The proof uses closed truth and falsehood formulas, so it also covers
zero-ary \(S\). The consequence he
draws, \(\mathrm{LFP}(\mathcal C)\subseteq\mathrm{PFP}(\mathcal C)\subseteq\mathrm{L}^\omega_{\infty\omega}(\mathcal C)\)
for every class \(\mathcal C\) of finite structures, is
\cref{thm:p03-w03-p03-059-lfp-in-pfp} together with
\cref{thm:p03-w03-p03-060-fp-in-lomega}, which also covers nested fixed
points. Kolaitis adds, without proof, that for positive \(\varphi\) every
transfinite stage is definable by an \(\mathrm{L}^k_{\infty\omega}\)-formula
and that LFP is contained in \(\mathrm{L}^\omega_{\infty\omega}\) on classes
of structures of bounded cardinality \cite[p.~54]{Kolaitis2003ExpressivePowerNotes};
neither statement is proved in this chapter.

\section{Datalog}
\label{sec:p03-w03-p03-060-datalog}

\begin{hbdefinition}{def:p03-w03-p03-060-datalog}{Datalog programs, Datalog\((\neq)\) and their stages}
Fix a vocabulary \(\sigma\) of input (extensional) symbols and new head
(intensional) relation symbols \(R_1,\ldots,R_n\) of arities
\(r_1,\ldots,r_n\ge0\). A \emph{rule} is an expression
\(R_i(\bar u)\leftarrow B_1\wedge\cdots\wedge B_m\) whose head is an atom of
a head symbol and whose body literals \(B_j\) are atoms \(S(\bar t)\) of
input or head symbols or equalities \(t=t'\), the terms being variables or
constant symbols of \(\sigma\). A \emph{Datalog program} \(\Pi\) is a finite
set of such rules; in a \emph{Datalog\((\neq)\) program} body literals may
also be inequalities \(t\neq t'\); Gr\"adel also admits negated input atoms
\(\neg S(\bar t)\), \(S\in\sigma\) \cite[Def.~3.3.57,
p.~180]{Gradel2007FMTDescriptiveChapter}. Head symbols never occur
negated. Kolaitis's Datalog has atoms only, and Datalog\((\neq)\) adds
inequalities \cite[pp.~42 and~44]{Kolaitis2003ExpressivePowerNotes}.

For a \(\sigma\)-structure \(\mathbf A\), the \emph{immediate-consequence
operator} \(\Gamma=(\Gamma_1,\ldots,\Gamma_n)\) maps
\(\bar X=(X_1,\ldots,X_n)\), \(X_l\subseteq A^{r_l}\), to the tuple whose
\(i\)-th entry is the set of \(\gamma(\bar u)\) for rules
\(R_i(\bar u)\leftarrow\beta\) of \(\Pi\) and maps \(\gamma\) from the
rule's variables into \(A\) under which \(\beta\) holds in
\((\mathbf A,\bar X)\). The \emph{stages} are
\(\bar X^0=(\varnothing,\ldots,\varnothing)\) and
\(\bar X^{m+1}=\Gamma(\bar X^m)\), and the query \((\Pi,R_i)\) assigns to
\(\mathbf A\) the \(i\)-th component of \(\bigcup_m\bar X^m\). This is
Gr\"adel's fixed-point semantics \cite[text before and in Def.~3.3.62,
p.~181]{Gradel2007FMTDescriptiveChapter}; its coincidence with the
minimal-model reading on finite databases is
\cref{sec:p09-presentation-n-c-024}.
\end{hbdefinition}

\begin{hblemma}{thm:p03-w03-p03-060-datalog-stages}{proved=annex:p03-w03-p03-060-datalog-stages}{The stages of a Datalog program reach the least fixed point by stage \(\omega\)}
For every Datalog\((\neq)\) program, possibly with negated input atoms, and
every \(\sigma\)-structure \(\mathbf A\), finite or infinite, \(\Gamma\) is
monotone in \(\bar X\), the stages increase, and
\(\bar X^\omega\defeq\bigcup_m\bar X^m\) is the least fixed point of
\(\Gamma\) and the least \(\bar Y\) with \(\Gamma(\bar Y)\subseteq\bar Y\)
(componentwise).
\end{hblemma}

Kolaitis notes that the closure ordinal of a Datalog\((\neq)\) program is
at most \(\omega\) on every infinite structure, because existential
quantification commutes with countable unions
\cite[p.~46]{Kolaitis2003ExpressivePowerNotes}.

\begin{hbtheorem}{thm:p03-w03-p03-060-datalog-in-existential}{proved=annex:p03-w03-p03-060-datalog-in-existential}{Datalog is contained in existential finite-variable infinitary logic}
Let \(\Pi\) be a Datalog program without negated input atoms and \(R_i\)
one of its head symbols. There are \(k\) and a formula
\(\varphi_i(x_1,\ldots,x_{r_i})\) of \(\exists\mathrm{L}_{\infty\omega}\)
with at most \(k\) distinct variables such that on every
\(\sigma\)-structure \(\mathbf A\), finite or infinite,
\((\Pi,R_i)(\mathbf A)=\{\bar a:\mathbf A\models\varphi_i[\bar x\mapsto\bar a]\}\).
For a Datalog\((\neq)\) program the same holds with
\(\exists\mathrm{L}_{\infty\omega}(\neq)\). One may take \(k\) to be the
largest number of variables in a rule plus \(2w\), where \(w\) is the
largest head arity.
\end{hbtheorem}

Combined with \cref{thm:p03-w03-p03-060-preservation}, the theorem yields
Kolaitis's preservation of Datalog queries under homomorphisms and of
Datalog\((\neq)\) queries under one-to-one homomorphisms
\cite[Props.~6.23 and~6.26, pp.~44--45]{Kolaitis2003ExpressivePowerNotes}.
The bound \(q+2w\) is not sharp. Kolaitis's sharp form keeps the number
\(k\) of variables of the rules, and it is a statement about systems of
existential formulas, of which Datalog programs are a special case.

\begin{hbdefinition}{def:p03-w03-p03-060-systems}{Systems of formulas and \(k\)-Datalog}
\begin{enumerate}
\item A \emph{system of formulas} over \(\sigma\) consists of distinct
      relation symbols \(S_1,\ldots,S_l\notin\sigma\) of arities
      \(n_1,\ldots,n_l\) and formulas
      \(\varphi_1(\bar x_1,S_1,\ldots,S_l)\), \dots,
      \(\varphi_l(\bar x_l,S_1,\ldots,S_l)\) of vocabulary \(\sigma\cup\{S_1,\ldots,S_l\}\), where \(\bar x_i\) is
      a tuple of \(n_i\) distinct variables containing the free variables of
      \(\varphi_i\). On a \(\sigma\)-structure \(\mathbf A\) it defines the
      operator
      \[
       \Phi_i(P_1,\ldots,P_l)=\{\bar a\in A^{n_i}:(\mathbf A,P_1,\ldots,P_l)\models\varphi_i[\bar x_i\mapsto\bar a]\},
      \]
      with stages \(\Phi^0=(\varnothing,\ldots,\varnothing)\) and
      \(\Phi^{m+1}=\Phi(\Phi^m)\). When \(\Phi\) is monotone, its least
      fixed point is written \((\varphi_1^\infty,\ldots,\varphi_l^\infty)\)
      \cite[Defs.~6.12 and~6.15, pp.~39--40]{Kolaitis2003ExpressivePowerNotes}.
\item A \emph{system of positive \(\exists\mathrm{FO}^k\)-formulas}
      (\emph{\(\exists\mathrm{FO}^k(\neq)\)-formulas}) is a system in which
      every \(\varphi_i\) belongs to
      \(\exists\mathrm{FO}^k(\sigma\cup\{S_1,\ldots,S_l\})\)
      (to \(\exists\mathrm{FO}^k(\neq)(\sigma\cup\{S_1,\ldots,S_l\})\)) and
      \(\bar x_i=(v_1,\ldots,v_{n_i})\); thus \(n_i\le k\).
\item A system is \emph{existential-conjunctive} (\emph{with
      inequalities}) when every \(\varphi_i\) is a finite disjunction of
      formulas \(\exists z_1\cdots\exists z_p\,\psi\), with
      \(z_1,\ldots,z_p\) distinct and \(\psi\) a finite conjunction of
      atomic formulas of \(\sigma\cup\{S_1,\ldots,S_l\}\), equalities
      included (and of inequalities \(t\neq t'\)).
\item A Datalog\((\neq)\) program is a \emph{\(k\)-Datalog\((\neq)\)
      program} when the body of every rule has at most \(k\) distinct
      variables and the head of every rule has at most \(k\) distinct
      variables; the body and the head may use different variables, so a
      rule may have up to \(2k\) variables \cite[p.~70]{Kolaitis2003ExpressivePowerNotes}.
\end{enumerate}
\end{hbdefinition}

\begin{hblemma}{thm:p03-w03-p03-060-continuity}{proved=annex:p03-w03-p03-060-continuity}{Existential formulas are continuous, and their systems close at \(\omega\)}
Let \(\varphi\) be a formula of \(\mathrm{L}_{\infty\omega}\) of vocabulary
\(\sigma\cup\{S_1,\ldots,S_l\}\) built from atomic formulas and
inequalities by \(\exists\), disjunctions of arbitrary sets and
conjunctions of finite sets.
\begin{enumerate}
\item \(\varphi\) is monotone in \(S_1,\ldots,S_l\).
\item If \(\bar P^0\subseteq\bar P^1\subseteq\cdots\) componentwise, then
      \((\mathbf A,\bigcup_n\bar P^n)\models\varphi[s]\) if and only if
      \((\mathbf A,\bar P^n)\models\varphi[s]\) for some \(n\).
\item For a system of such formulas, on every \(\sigma\)-structure, finite
      or infinite, the stages increase and \(\bigcup_m\Phi^m\) is the least
      fixed point of \(\Phi\).
\end{enumerate}
\end{hblemma}

Item~3 is Kolaitis's observation that the closure ordinal of a
Datalog\((\neq)\) program is at most \(\omega\) on every infinite structure
\cite[p.~46]{Kolaitis2003ExpressivePowerNotes}, extended to the
systems used below.

\begin{hbtheorem}{thm:p03-w03-p03-060-kv-systems}{proved=annex:p03-w03-p03-060-kv-systems}{Kolaitis and Vardi: systems of positive \(\exists\mathrm{FO}^k\)-formulas}
Let \(k\ge1\) and let
\[
 \varphi_1(v_1,\ldots,v_{n_1},S_1,\ldots,S_l),\ \ldots,\ \varphi_l(v_1,\ldots,v_{n_l},S_1,\ldots,S_l)
\]
be a system of positive \(\exists\mathrm{FO}^k\)-formulas over \(\sigma\),
with operator \(\Phi\).
\begin{enumerate}
\item For every \(m\in\N\) and \(i\le l\) there is a formula
      \(\varphi^m_i\in\exists\mathrm{FO}^k(\sigma)\) with free labels in
      \([n_i]\) that defines the component \(\Phi^m_i\) of the stage
      \(\Phi^m\) on the class of all \(\sigma\)-structures.
\item The formula \(\bigvee_{m\in\N}\varphi^m_i\) of
      \(\exists\mathrm{L}^k_{\infty\omega}(\sigma)\) defines the component
      \(\varphi_i^\infty\) of the least fixed point on the class of all
      \(\sigma\)-structures.
\end{enumerate}
The same holds with \(\exists\mathrm{FO}^k(\neq)\) and
\(\exists\mathrm{L}^k_{\infty\omega}(\neq)\).
\end{hbtheorem}

This is Kolaitis's Theorem~8.2, for which he gives a hint and a reference
to Kolaitis and Vardi \cite[Thm.~8.2, p.~70]{Kolaitis2003ExpressivePowerNotes}.
His second item, as printed, says that each component of the least fixed
point is \(\exists\mathrm{FO}^k\)-definable; the consequence he draws in the
same statement is \(\exists\mathrm{L}^k_{\infty\omega}\)-definability, and
the first-order reading cannot hold
(\cref{thm:p03-w03-p03-060-tc-not-fo}). Item~2 above is the infinitary
statement.

\begin{hbtheorem}{thm:p03-w03-p03-060-kdatalog-system}{proved=annex:p03-w03-p03-060-kdatalog-system}{\(k\)-Datalog programs as systems of positive \(\exists\mathrm{FO}^k\)-formulas}
Let \(\Pi\) be a \(k\)-Datalog program (a \(k\)-Datalog\((\neq)\) program)
with head symbols \(R_1,\ldots,R_n\) of arities \(r_1,\ldots,r_n\) and
stages \(\bar X^m\) (\cref{def:p03-w03-p03-060-datalog}). For \(i\le n\) let
\(M_i\) be the finite set of maps \(\mu\colon[r_i]\to T_k\), and for
\(\alpha\colon[k]\to A\) write
\(\mu^\alpha\defeq(\mu(1)^{\mathbf A}[\alpha],\ldots,\mu(r_i)^{\mathbf A}[\alpha])\).
There is a system of positive \(\exists\mathrm{FO}^k\)-formulas
(\(\exists\mathrm{FO}^k(\neq)\)-formulas)
\(\varphi_{i,\mu}(v_1,\ldots,v_k,\bar S)\), \(i\le n\), \(\mu\in M_i\), over
symbols \(S_{i,\mu}\) of arity \(k\), such that on every
\(\sigma\)-structure, for every \(m\), \(i\) and \(\mu\),
\[
 \Phi^m_{i,\mu}=\{\alpha\in A^k:\mu^\alpha\in X^m_i\},\qquad
 \varphi^\infty_{i,\mu}=\{\alpha\in A^k:\mu^\alpha\in(\Pi,R_i)(\mathbf A)\},
\]
where \(\alpha\in A^k\) is read as the assignment \(j\mapsto\alpha(j)\).
\end{hbtheorem}

Kolaitis asserts this as a refinement of the proof of his
Proposition~6.21 \cite[p.~71]{Kolaitis2003ExpressivePowerNotes}. A head
symbol may have arity larger than \(k\), since a head with at most \(k\)
distinct variables may repeat them; the maps \(\mu\) record which
\(k\)-variable pattern a tuple of the relation is read through, and this
keeps every formula of the system inside \(\exists\mathrm{FO}^k\). The
bound on the variables of the head is not used.

\begin{hbcorollary}{thm:p03-w03-p03-060-kdatalog}{proved=annex:p03-w03-p03-060-kdatalog}{\(k\)-Datalog is contained in \(\exists\mathrm{L}^k_{\infty\omega}\), and \(k\)-Datalog\((\neq)\) in \(\exists\mathrm{L}^k_{\infty\omega}(\neq)\)}
Let \(\Pi\) be a \(k\)-Datalog program and \(R_i\) a head symbol of arity
\(r_i\le k\). Every stage \(X^m_i\) is defined on all \(\sigma\)-structures
by a formula of \(\exists\mathrm{FO}^k\) with free labels in \([r_i]\), and
there is \(\chi_i\in\exists\mathrm{L}^k_{\infty\omega}(\sigma)\) with free
labels in \([r_i]\) such that for every \(\sigma\)-structure \(\mathbf A\),
finite or infinite, and every \(\alpha\colon[k]\to A\),
\(\mathbf A\models\chi_i[\alpha]\iff(\alpha(1),\ldots,\alpha(r_i))\in(\Pi,R_i)(\mathbf A)\).
In symbols, \(k\text{-Datalog}(\mathcal S)\subseteq\exists\mathrm{L}^k_{\infty\omega}(\mathcal S)\)
for the class \(\mathcal S\) of all \(\sigma\)-structures. The same holds
for \(k\)-Datalog\((\neq)\), \(\exists\mathrm{FO}^k(\neq)\) and
\(\exists\mathrm{L}^k_{\infty\omega}(\neq)\).
\end{hbcorollary}

These are the conclusion of Kolaitis's Theorem~8.2 and his Theorem~8.3
\cite[pp.~70--71]{Kolaitis2003ExpressivePowerNotes}. The arity
hypothesis \(r_i\le k\) is the reading under which the inclusion is
meaningful: a query of arity \(r_i\) is defined by a formula with free
labels \(1,\ldots,r_i\), which needs \(r_i\) variables. Every Datalog
program is a \(k\)-Datalog program for \(k\) the largest number of
variables in a rule body or head, so the corollary sharpens the bound of
\cref{thm:p03-w03-p03-060-datalog-in-existential}. Kolaitis's display of a
path of length three in \(\exists\mathrm{FO}^3\) is the case of the
transitive-closure program, which is a \(3\)-Datalog program
\cite[p.~70]{Kolaitis2003ExpressivePowerNotes}.

\begin{hblemma}{thm:p03-w03-p03-060-tc-not-fo}{proved=annex:p03-w03-p03-060-tc-not-fo}{Least fixed points of \(\exists\mathrm{FO}^k\)-systems need infinitary disjunctions}
No first-order formula \(\theta(x,y)\) of vocabulary \(\{E\}\), with any
number of variables, defines on the class of all \(\{E\}\)-structures the
set of pairs \((a,b)\) joined by a directed path of length at least
\(1\). Hence, although the transitive closure is the least fixed point of
the system consisting of the single \(\exists\mathrm{FO}^3\)-formula
\(E(v_1,v_2)\vee\exists v_3(E(v_1,v_3)\wedge S(v_3,v_2))\), it is not
\(\exists\mathrm{FO}^3\)-definable on that class.
\end{hblemma}

The other bridge between Datalog and fixed points runs in both directions
and does not count variables. Kolaitis attributes it to Chandra and Harel.

\begin{hbtheorem}{thm:p03-w03-p03-060-datalog-systems}{proved=annex:p03-w03-p03-060-datalog-systems}{Datalog programs are existential-conjunctive systems}
\begin{enumerate}
\item For every Datalog program (Datalog\((\neq)\) program) \(\Pi\) with
      head symbols \(R_1,\ldots,R_n\) there is an existential-conjunctive
      system (with inequalities) over the symbols \(R_1,\ldots,R_n\) whose
      operator is \(\Gamma\) on every \(\sigma\)-structure; hence
      \(\varphi_i^\infty=(\Pi,R_i)(\mathbf A)\) for every \(\mathbf A\).
\item For every existential-conjunctive system (with inequalities)
      \(\varphi_1(\bar x_1),\ldots,\varphi_l(\bar x_l)\) over
      \(S_1,\ldots,S_l\) there is a Datalog (Datalog\((\neq)\)) program with
      head symbols \(S_1,\ldots,S_l\) whose operator \(\Gamma\) is the
      operator of the system on every \(\sigma\)-structure; hence
      \((\Pi,S_i)(\mathbf A)=\varphi_i^\infty\) for every \(\mathbf A\).
\item If \(\sigma\) has no constant symbols, every Datalog\((\neq)\)
      program has the same operator \(\Gamma\) as a program with the same
      head symbols whose rule bodies contain no equality literal.
\end{enumerate}
Consequently, for every class \(\mathcal C\) of \(\sigma\)-structures, a
query on \(\mathcal C\) is definable by a Datalog program if and only if it
is defined on \(\mathcal C\) by \(\varphi_m^\infty\) for some
existential-conjunctive system \(\varphi_1,\ldots,\varphi_m\), and it is
definable by a Datalog\((\neq)\) program if and only if this holds for some
existential-conjunctive system with inequalities.
\end{hbtheorem}

The two consequences are Kolaitis's Proposition~6.21, which he credits to
Chandra and Harel, and the analogue for Datalog\((\neq)\) that he states
after Example~6.24 \cite[pp.~43 and~45]{Kolaitis2003ExpressivePowerNotes}.
His rules contain relational atoms only, with no equality literals
\cite[p.~42]{Kolaitis2003ExpressivePowerNotes}; item~3 reconciles this
with the equalities that the systems may contain, for vocabularies without
constants. His hint reads one rule as one formula and one formula as one
rule \cite[p.~44]{Kolaitis2003ExpressivePowerNotes}. The statement in
his text requires each \(\varphi_i\) to be a single formula
\(\exists\bar z_i\,\psi_i\); a program in which some head symbol has two
rules needs the finite disjunction of item~3 of
\cref{def:p03-w03-p03-060-systems}, as his own formulas for Examples~6.20
and~6.24 do \cite[pp.~43 and~45]{Kolaitis2003ExpressivePowerNotes}.
The disjunction cannot be dropped.

\begin{hblemma}{thm:p03-w03-p03-060-single-conjunction}{proved=annex:p03-w03-p03-060-single-conjunction}{Single conjunctions do not suffice}
For \(\sigma\)-structures \(\mathbf A,\mathbf B\) let
\(\mathbf A\times\mathbf B\) have universe \(A\times B\), \(R^{\mathbf A\times\mathbf B}\) the set of
tuples \(((a_1,b_1),\ldots,(a_r,b_r))\) with \(\bar a\in R^{\mathbf A}\)
and \(\bar b\in R^{\mathbf B}\), and
\(c^{\mathbf A\times\mathbf B}=(c^{\mathbf A},c^{\mathbf B})\); write
\(\langle\bar a,\bar b\rangle=((a_1,b_1),\ldots,(a_n,b_n))\).
\begin{enumerate}
\item Let \(\varphi_1,\ldots,\varphi_l\) be a system in which every
      \(\varphi_i\) is a single formula \(\exists\bar z\,\psi\) with
      \(\psi\) a finite conjunction of atomic formulas and inequalities.
      If \(\bar a\in\varphi_i^\infty(\mathbf A)\) and
      \(\bar b\in\varphi_i^\infty(\mathbf B)\), then
      \(\langle\bar a,\bar b\rangle\in\varphi_i^\infty(\mathbf A\times\mathbf B)\).
\item The transitive closure of \(E\) does not have this property, so it
      is \(\varphi_i^\infty\) for no such system, on the class of all
      finite graphs; it is defined by the two-rule Datalog program
      \(S(x,y)\leftarrow E(x,y)\); \(S(x,y)\leftarrow E(x,z)\wedge S(z,y)\)
      \cite[Ex.~6.18, p.~42]{Kolaitis2003ExpressivePowerNotes}.
\end{enumerate}
\end{hblemma}

Gr\"adel's form of the correspondence compares Datalog with a fragment of
fixed-point logic in which simultaneous fixed points are allowed.

\begin{hbdefinition}{def:p03-w03-p03-060-efp}{Existential fixed-point logic EFP}
Let \(\tau\) be a vocabulary (free relation variables are treated as its
symbols). The formulas of \(\mathrm{EFP}(\tau)\) are generated by:
\begin{enumerate}
\item every \emph{literal} \(P(\bar t)\), \(\neg P(\bar t)\), \(t=t'\),
      \(\neg\,t=t'\) with \(P\in\tau\) and terms \(\bar t,t,t'\);
\item \(\psi_1\wedge\psi_2\), \(\psi_1\vee\psi_2\) and \(\exists x\,\psi\);
\item \([\mathbf{lfp}\,R_j:\mathcal S](\bar t)\), where \(\mathcal S\) is a
      system \(R_1\bar x_1\defeq\varphi_1,\ldots,R_m\bar x_m\defeq\varphi_m\)
      of new relation symbols \(R_l\notin\tau\) of arities \(k_l\),
      \(k_l\)-tuples \(\bar x_l\) of distinct variables and formulas
      \(\varphi_l\in\mathrm{EFP}(\tau\cup\{R_1,\ldots,R_m\})\) containing
      no literal \(\neg R_n(\bar s)\), \(j\le m\), and \(\bar t\) is a
      \(k_j\)-tuple of terms.
\end{enumerate}
The \emph{parameters} of \(\mathcal S\) are the variables free in some
\(\varphi_l\) outside \(\bar x_l\); the free variables of
\([\mathbf{lfp}\,R_j:\mathcal S](\bar t)\) are those of \(\bar t\) and the
parameters. In a \(\tau\)-structure, with the parameters assigned, the
system has the operator
\(F_{\mathcal S}(\bar Y)_l=\{\bar a:(\mathbf A,\bar Y)\models\varphi_l[\bar x_l\mapsto\bar a]\}\),
which is monotone (annex to \cref{thm:p03-w03-p03-060-efp-datalog}, Step~0),
and \([\mathbf{lfp}\,R_j:\mathcal S](\bar t)\) holds when the value of
\(\bar t\) lies in the \(j\)-th component of
\(\operatorname{lfp}(F_{\mathcal S})\), which exists by
\cref{thm:p03-w03-p03-059-knaster-tarski} applied to the power set of the
disjoint union of the \(A^{k_l}\). These are Gr\"adel's simultaneous
fixed points \cite[Def.~3.3.39, p.~171]{Gradel2007FMTDescriptiveChapter}
restricted as in his definition of EFP: no universal quantifier, no
\(\mathbf{gfp}\), negation on atomic formulas only \cite[Def.~3.3.63,
p.~182]{Gradel2007FMTDescriptiveChapter}.
\end{hbdefinition}

\begin{hbtheorem}{thm:p03-w03-p03-060-efp-datalog}{proved=annex:p03-w03-p03-060-efp-datalog}{Gr\"adel: Datalog is equivalent to EFP}
Datalog here is Gr\"adel's: Datalog\((\neq)\) programs whose rule bodies
may also contain negated input atoms.
\begin{enumerate}
\item For every Datalog program \(\Pi\) and head symbol \(R_i\) there is an
      EFP formula \([\mathbf{lfp}\,R_i:\mathcal S_\Pi](x_1,\ldots,x_{r_i})\)
      that defines \((\Pi,R_i)\) on every structure, finite or infinite.
\item For every \(\psi\in\mathrm{EFP}(\tau)\) and every tuple
      \(\bar z=(z_1,\ldots,z_n)\) of distinct variables containing the free
      variables of \(\psi\) there are a Datalog program \(\Pi_\psi\) with
      input symbols in \(\tau\) and a head symbol \(H_\psi\) of arity \(n\)
      such that for every \(\tau\)-structure \(\mathbf A\), finite or
      infinite,
      \((\Pi_\psi,H_\psi)(\mathbf A)=\{\bar c\in A^n:\mathbf A\models\psi[\bar z\mapsto\bar c]\}\).
\item In both translations a negated input atom \(\neg P(\bar t)\) occurs
      on one side only if a negated \(P\)-literal occurs on the other, and
      an inequality only if a negated equality does.
\end{enumerate}
Hence Datalog and EFP define the same queries; Datalog\((\neq)\) and the
EFP formulas in which negation stands only before equalities define the
same queries; and Datalog and the EFP formulas without negation define the
same queries.
\end{hbtheorem}

This is Gr\"adel's Proposition~3.3.64, stated for finite databases, with
the inclusion of Datalog in EFP derived in the text before it and the
converse left as an exercise \cite[pp.~181--182]{Gradel2007FMTDescriptiveChapter}.
The annex proves the converse on all structures. Its one delicate point is
the parameters of a system: the relations \(R_l\) of an inner fixed point
depend on the values of the parameters, while a head relation of a program
has one value per structure. The translation therefore adds the parameters
as extra arguments of each \(R_l\).

\begin{hbexample}{def:p03-w03-p03-060-efp-example}{A fixed point with a parameter as a program}
Let \(\psi(y,p)\defeq[\mathbf{lfp}\,R\,x\,.\,x=p\vee\exists z(R(z)\wedge E(z,x))](y)\),
with the parameter \(p\): \(y\) is reachable from \(p\) by a directed walk
of length at least \(0\). Up to the auxiliary head symbols that copy one
relation into another, the program of
\cref{thm:p03-w03-p03-060-efp-datalog} is
\[
 R^+(x,p)\leftarrow x=p;\qquad
 R^+(x,p)\leftarrow R^+(z,p)\wedge E(z,x);\qquad
 H(y,p)\leftarrow R^+(y,p),
\]
in which \(R^+\) holds of \((a,c)\) when \(a\) is reachable from \(c\).
Dropping the extra argument gives the rule \(R(x)\leftarrow x=p\), in which
\(p\) is a variable of the rule and ranges over all elements, so \(R\)
would contain every element: the parameter would be lost.
\end{hbexample}

\begin{hbtheorem}{thm:p03-w03-p03-060-datalog-in-lfp}{proved=annex:p03-w03-p03-060-datalog-in-lfp}{Datalog is a fragment of LFP}
Let \(\Pi\) be a Datalog\((\neq)\) program, possibly with negated input
atoms, and \(R_i\) a head symbol. There is an LFP formula
\(\chi_i(x_1,\ldots,x_{r_i})\) of vocabulary \(\sigma\), using only
\(\mathbf{lfp}\) constructors (no \(\mathbf{gfp}\)) applied to formulas
without universal quantifiers in which negation is applied only to input
atoms and equalities, such that on every finite \(\sigma\)-structure
\(\chi_i\) defines \((\Pi,R_i)\).
\end{hbtheorem}

Gr\"adel obtains the query as a simultaneous least fixed point of
existential update formulas \cite[pp.~181--182]{Gradel2007FMTDescriptiveChapter},
calls the resulting fragment existential fixed-point logic EFP
\cite[Def.~3.3.63, p.~182]{Gradel2007FMTDescriptiveChapter}, and states
the equivalence of Datalog with EFP with the converse left as an exercise
\cite[Prop.~3.3.64, p.~182]{Gradel2007FMTDescriptiveChapter}; simultaneous
inductions are removed by nesting \cite[Lemma~3.3.41 and Thm.~3.3.42,
pp.~172--173]{Gradel2007FMTDescriptiveChapter}. The annex proves the
inclusion with single, nested fixed points directly; the converse
inclusion of EFP in Datalog is \cref{thm:p03-w03-p03-060-efp-datalog}.

\begin{hbexample}{def:p03-w03-p03-060-tc-example}{Transitive closure, step by step}
Take \(\sigma=\{E\}\) and the program
\(S(x,y)\leftarrow E(x,y)\); \(S(x,y)\leftarrow E(x,z)\wedge S(z,y)\)
\cite[Ex.~6.18, p.~42]{Kolaitis2003ExpressivePowerNotes}. Its stage
\(m\) consists of the pairs joined by a directed path of length between
\(1\) and \(m\). The construction of
\cref{thm:p03-w03-p03-060-datalog-in-existential} produces, with head
variables \(x_1,x_2\) and carriers \(z_1,z_2\), formulas \(\varphi^m\) whose
disjunction over \(m\) defines \(S\) with the seven variables
\(x,y,z,x_1,x_2,z_1,z_2\), the bound of the theorem for three rule
variables and head arity \(2\). Reusing variables by hand does
better: the formulas
\(\pi_1(x,y)\defeq E(x,y)\) and
\(\pi_{m+1}(x,y)\defeq\exists z\bigl(E(x,z)\wedge\exists x(x=z\wedge\pi_m(x,y))\bigr)\)
lie in \(\exists\mathrm{FO}^3\), and \(\bigvee_{m\ge1}\pi_m\) defines \(S\)
on all graphs; Kolaitis displays the case \(m=3\)
\cite[p.~70]{Kolaitis2003ExpressivePowerNotes}, and
\cref{thm:p03-w03-p03-056-walks} gives the general three-variable
construction. The formula given by
\cref{thm:p03-w03-p03-060-datalog-in-lfp} is equivalent on finite graphs to
\([\mathbf{lfp}\,S x_1x_2\,.\,E(x_1,x_2)\vee\exists z(E(x_1,z)\wedge S(z,x_2))](x,y)\),
which is Gr\"adel's Example~3.3.8 up to the names of the
variables \cite[p.~155]{Gradel2007FMTDescriptiveChapter}.
\end{hbexample}

\section{Comparing the fixed-point logics}
\label{sec:p03-w03-p03-060-comparison}

\begin{hbtheorem}{thm:p03-w03-p03-060-fp-comparison}{proved=annex:p03-w03-p03-060-fp-comparison}{Least, inflationary and partial fixed points on finite structures}
Let \(\mathbf A\) be a finite structure, \(\psi(R,\bar x)\) a formula of
LFP, IFP or PFP and \(\bar t\) a tuple of terms of the arity of \(R\); fix
the parameters.
\begin{enumerate}
\item If \(F_\psi\) is monotone (in particular, by
      \cref{thm:p03-w03-p03-059-positive-monotone}, if \(R\) occurs only
      positively in an LFP formula \(\psi\)), then
      \([\mathbf{lfp}\,R\bar x\,.\,\psi](\bar t)\Leftrightarrow
      [\mathbf{ifp}\,R\bar x\,.\,\psi](\bar t)\).
\item For every \(\psi\),
      \([\mathbf{ifp}\,R\bar x\,.\,\psi](\bar t)\Leftrightarrow
      [\mathbf{pfp}\,R\bar x\,.\,R(\bar x)\vee\psi](\bar t)\).
\end{enumerate}
Consequently every LFP formula is equivalent on finite structures to an
IFP formula, and every IFP formula to a PFP formula:
\(\mathrm{LFP}\le\mathrm{IFP}\le\mathrm{PFP}\) on finite structures.
\end{hbtheorem}

The first item is the finite case of Gr\"adel's Exercise~3.3.44(4)
\cite[p.~174]{Gradel2007FMTDescriptiveChapter}; the direct inclusion of LFP
in PFP is \cref{thm:p03-w03-p03-059-lfp-in-pfp}. The reverse inclusions are
of a different order of difficulty: IFP into LFP uses the stage comparison
below, while \cref{thm:p03-w03-p03-059-abiteboul-vianu} identifies
PFP into LFP on all finite structures with
\(\mathsf P=\mathsf{PSPACE}\).

The converse of the first inclusion rests on comparing the stages at which
tuples enter an inflationary induction.

\begin{hbdefinition}{def:p03-w03-p03-060-stage-comparison}{Inflationary rank and stage comparison}
Let \(\mathbf A\) be a finite structure, \(\varphi(X,\bar x)\) a formula
with \(X\) of arity \(r\ge1\), the parameters fixed, and \(X^m\) the
inflationary stages of \cref{def:p03-w03-p03-060-ifp}, with inflationary
fixed point \(X^\infty\). The \emph{rank} of \(\bar a\in A^r\) is
\(|\bar a|_\varphi\defeq\min\{m:\bar a\in X^m\}\) if
\(\bar a\in X^\infty\) and \(|\bar a|_\varphi\defeq\infty\) otherwise. The
\emph{inflationary stage comparison relations} are
\[
 \bar a\le^{\mathrm{inf}}_\varphi\bar b\iff|\bar a|_\varphi\le|\bar b|_\varphi<\infty,
 \qquad
 \bar a\prec^{\mathrm{inf}}_\varphi\bar b\iff|\bar a|_\varphi<|\bar b|_\varphi .
\]
These are Gr\"adel's relations for the inflationary operator
\(S\mapsto S\cup F_\varphi(S)\) \cite[Def.~3.3.47,
p.~175]{Gradel2007FMTDescriptiveChapter}. Since stage \(0\) is empty, every
finite rank is at least \(1\); \(\bar a\in X^\infty\) exactly when
\(\bar a\le^{\mathrm{inf}}_\varphi\bar a\).
\end{hbdefinition}

\begin{hblemma}{thm:p03-w03-p03-060-stage-comparison}{proved=annex:p03-w03-p03-060-stage-comparison}{Inflationary stage comparison is LFP-definable on finite structures}
Let \(\varphi(X,\bar x)\) be an LFP formula of vocabulary
\(\tau\cup\{X\}\) in which the relation variable \(X\), of arity \(r\ge1\),
may occur both positively and negatively; let \(\bar x=(x_1,\ldots,x_r)\)
be distinct variables and \(\bar p\) the free first-order variables of
\(\varphi\) other than those of \(\bar x\). Let \(\bar u,\bar w\) be
\(r\)-tuples of variables, all \(2r\) distinct, none occurring in
\(\varphi\) and none among \(\bar x\). There are LFP formulas \(\lambda_\varphi(\bar u,\bar w)\),
\(\mathrm{le}_\varphi(\bar u,\bar w)\) and \(\mathrm{lt}_\varphi(\bar u,\bar w)\)
of vocabulary \(\tau\), with free first-order variables among
\(\bar u,\bar w,\bar p\), such that on every finite \(\tau\)-structure and
for every assignment to \(\bar p\) they define, respectively,
\(\{(\bar a,\bar b):|\bar a|_\varphi\le|\bar b|_\varphi,\ |\bar a|_\varphi<\infty\}\),
\(\le^{\mathrm{inf}}_\varphi\) and \(\prec^{\mathrm{inf}}_\varphi\). The
free relation variables of \(\varphi\) other than \(X\) are counted in
\(\tau\) and may have any interpretation.
\end{hblemma}

This is the finite case of Gr\"adel's inflationary stage comparison
theorem, which he states for formulas of FO or LFP and for which he refers
to the literature for proofs \cite[Thm.~3.3.49,
p.~176]{Gradel2007FMTDescriptiveChapter}. Gr\"adel adds that on finite
structures the relations are even definable in positive LFP. The formulas
built in the annex may contain negated fixed points; Immerman's normal form
removes them.

\begin{hbcorollary}{thm:p03-w03-p03-060-stage-comparison-positive}{proved=annex:p03-w03-p03-060-stage-comparison-positive}{Inflationary stage comparison is definable in positive LFP on finite structures}
Under the hypotheses of \cref{thm:p03-w03-p03-060-stage-comparison} there
are formulas \(\mathrm{le}^+_\varphi(\bar u,\bar w)\) and
\(\mathrm{lt}^+_\varphi(\bar u,\bar w)\) of \(\mathrm{LFP}_1\)
(\cref{def:p03-w03-p03-059-lfp1}) of vocabulary \(\tau\), with free
first-order variables among \(\bar u,\bar w,\bar p\), that define
\(\le^{\mathrm{inf}}_\varphi\) and \(\prec^{\mathrm{inf}}_\varphi\) on every
finite \(\tau\)-structure, for every assignment to \(\bar p\); the
vocabulary \(\tau\) may contain constant symbols.
\end{hbcorollary}

This is the second sentence of Gr\"adel's Theorem~3.3.49, for \(\varphi\)
in FO or LFP \cite[Thm.~3.3.49, p.~176]{Gradel2007FMTDescriptiveChapter}.
The proof combines \cref{thm:p03-w03-p03-060-stage-comparison} with
\cref{thm:p03-w03-p03-059-immerman-normal-form}, which is stated for
relational vocabularies; constant symbols are first replaced by
parameters.

\begin{hbtheorem}{thm:p03-w03-p03-060-gurevich-shelah}{proved=annex:p03-w03-p03-060-gurevich-shelah}{Gurevich and Shelah: IFP collapses to LFP on finite structures}
For every IFP formula \(\varphi\) of vocabulary \(\sigma\) there is an LFP
formula \(\varphi^\dagger\) of vocabulary \(\sigma\), with the same free
relation variables and with free first-order variables among those of
\(\varphi\), such that \(\varphi\) and \(\varphi^\dagger\) have the same
truth value in every finite structure, under every assignment and every
interpretation of the free relation variables. Together with
\cref{thm:p03-w03-p03-060-fp-comparison}, LFP and IFP have the same
expressive power on finite structures.
\end{hbtheorem}

Gr\"adel attributes the finite case to Gurevich and Shelah and notes that
their proof does not carry over to infinite structures
\cite[p.~175]{Gradel2007FMTDescriptiveChapter}; Kolaitis records the same
equivalence on classes of finite structures
\cite[p.~49]{Kolaitis2003ExpressivePowerNotes}. The annex proof
follows the route Gr\"adel indicates, through a stage comparison lemma
\cite[pp.~175--176]{Gradel2007FMTDescriptiveChapter}; the arity of the
fixed-point variables doubles, as Gr\"adel warns that it must in general
\cite[p.~176]{Gradel2007FMTDescriptiveChapter}.

\begin{hbtheorem}{thm:p03-w03-p03-060-kreutzer}{cited=Gradel2007FMTDescriptiveChapter}{Kreutzer: IFP collapses to LFP on all structures}
When the inflationary stages are continued through the ordinals, with
unions at limit stages, every IFP formula is equivalent on all structures,
finite or infinite, to an LFP formula.
\end{hbtheorem}

This is Gr\"adel's Theorem~3.3.50, derived there in one line from the
general stage comparison Theorem~3.3.49, whose proof he does not give
\cite[p.~176]{Gradel2007FMTDescriptiveChapter}. The handbook proves only the
finite case, \cref{thm:p03-w03-p03-060-gurevich-shelah}; on infinite
structures the translation also introduces alternations of least and
greatest fixed points \cite[p.~176]{Gradel2007FMTDescriptiveChapter}.

The comparison of LFP with PFP on all finite structures is the theorem of
Abiteboul and Vianu, \cref{thm:p03-w03-p03-059-abiteboul-vianu}: the two
logics are equivalent on finite structures if and only if
\(\mathsf P=\mathsf{PSPACE}\) \cite[Thm.~3.3.56,
p.~180]{Gradel2007FMTDescriptiveChapter}, where Gr\"adel states it without
proof. The following corollary isolates the implication to equality
of the complexity classes. The full converse is proved through
\cref{sec:p03-abiteboul-vianu-completion}.

\begin{hbcorollary}{thm:p03-w03-p03-060-av-easy}{proved=annex:p03-w03-p03-060-av-easy}{Equivalence of the fixed-point logics forces \(\mathsf P=\mathsf{PSPACE}\)}
Consider the two statements:
\begin{enumerate}
\item[(E\(_{\mathrm{LFP}}\))] for every vocabulary \(\sigma\) and every PFP
      formula \(\psi\) of vocabulary \(\sigma\) there is an LFP formula of
      vocabulary \(\sigma\), with free variables among those of \(\psi\),
      equivalent to \(\psi\) on all finite \(\sigma\)-structures;
\item[(E\(_{\mathrm{IFP}}\))] the same with an IFP formula in place of the
      LFP formula.
\end{enumerate}
Then (E\(_{\mathrm{LFP}}\)) and (E\(_{\mathrm{IFP}}\)) are equivalent, and
each implies \(\mathsf P=\mathsf{PSPACE}\).
\end{hbcorollary}

For the converse, \cref{thm:p03-av-type-order} constructs the
LFP-definable order of tuple types. The quotient interpretation and
its polynomial-space evaluation are
\cref{thm:p03-av-pullback,thm:p03-av-quotient-space}. Their combination
with ordered capture proves (E\(_{\mathrm{LFP}}\)) from
\(\mathsf P=\mathsf{PSPACE}\), including nested fixed points and
captured parameters. Kolaitis outlines the ordered-type route in
\cite[Thms.~7.24 and~7.25, p.~68]{Kolaitis2003ExpressivePowerNotes}.

\subsection{The unordered fixed-point equivalence}
\label{sec:p03-abiteboul-vianu-completion}

The construction orders tuple types, not the original universe.
Its ordered structure retains the information needed by a fixed
formula without choosing an order on the original elements.

The type-invariant method of Dawar, Lindell and Weinstein is presented
in \cite[Chapter~11]{Libkin2004}. The evaluation below also retains
nested fixed-point binders and their parameters.

\begin{hbconvention}{def:p03-av-setting}{Finite structures and tuple types}
Use the finite vocabularies of
\cref{def:p03-w03-p03-056-vocabulary}: relation symbols of positive
arity and possibly constants, with finite nonempty universes.
For \(k\geq1\), two tuples \(\bar a,\bar b\in A^k\)
have the same \(k\)-variable type when they satisfy the same
first-order formulas using the variables \(x_1,\ldots,x_k\).
This is \(\equiv^k_{\omega\omega}\) of
\cref{def:p03-w03-p03-056-equivalence}, with the variables renamed.
Write \(\bar a[i\leftarrow u]\) for replacement of coordinate
\(i\). This is notation for a tuple, not a function symbol added
to the object language.
\end{hbconvention}

\begin{hblemma}{thm:p03-av-type-order}{proved=annex:p03.av.1}{Ordering the tuple types}
For each vocabulary and \(k\), there is an LFP formula
\(D_k(\bar x,\bar y)\) such that, in every finite structure,
\[
 E_k(\bar x,\bar y)
 \quad\defeq\quad
 \neg D_k(\bar x,\bar y)\land\neg D_k(\bar y,\bar x)
\]
defines equality of \(k\)-variable types, and \(D_k\) induces a
strict linear order on the \(E_k\)-classes. The formula is uniform
in the input structure. It need not define a linear order on \(A\).
\end{hblemma}

\begin{hbdefinition}{def:p03-av-quotient}{The ordered type structure}
Fix the formula \(D_k\) just constructed. The universe of
\(\operatorname{Typ}_k(\mathbf A)\) is \(A^k/E_k\). Its vocabulary
is finite and depends only on \(k\) and the original vocabulary.
It contains the induced strict order and a constant \(c_\ast\)
naming its least element, together with the following relations.
\begin{enumerate}
\item For each atomic formula \(\alpha\) in the fixed variables
      and constants, a unary predicate \(U_\alpha\), true of
      \([\bar a]\) exactly when \(\mathbf A\models\alpha[\bar a]\).
\item For each coordinate \(i\), a binary relation \(S_i\), where
      \(S_i([\bar a],[\bar b])\) means that some representatives
      agree outside coordinate \(i\). Equivalently,
      \(\exists u\,E_k(\bar a[i\leftarrow u],\bar b)\).
\item Each coordinate map \(\theta\) selects, at every position,
      one of the \(k\) fixed variables or an original constant.
      Include a binary relation \(P_\theta\) with graph
      \[
       P_\theta([\bar a],[\bar b])
             \quad\Longleftrightarrow\quad E_k(\bar a_\theta,\bar b).
      \]
      The \(i\)-th entry of \(\bar a_\theta\) is the value of
      \(\theta(i)\) under \(\bar a\).
\end{enumerate}
There are finitely many atomic formulas and maps in these lists.
The next lemma verifies that the definitions do not depend on
representatives.
\end{hbdefinition}

\begin{hblemma}{thm:p03-av-pullback}{proved=annex:p03.av.2}{Interpreting the ordered quotient}
\(\operatorname{Typ}_k(\mathbf A)\) is well defined and has between
one and \(|A|^k\) elements. Its relations and its quotient equality
have uniform LFP definitions on representative tuples. Every LFP
sentence \(\eta\) in its vocabulary has an LFP sentence
\(\eta^{\leftarrow}\) in the original vocabulary satisfying
\[
 \operatorname{Typ}_k(\mathbf A)\models\eta
 \quad\Longleftrightarrow\quad
 \mathbf A\models\eta^{\leftarrow}.
\]
The construction handles \(c_\ast\) by its definable equivalence
class; it does not choose a definable representative from that class.
\end{hblemma}

\begin{hblemma}{thm:p03-av-quotient-space}{proved=annex:p03.av.3}{Evaluation on types}
For every PFP sentence \(\varphi\) there are a fixed \(k\)
and an isomorphism-invariant polynomial-space decision problem
\(Q_\varphi\) on finite ordered structures of the vocabulary
of \(\operatorname{Typ}_k\), with its distinguished constant, such that
\[
 \mathbf A\models\varphi
 \quad\Longleftrightarrow\quad
 \operatorname{Typ}_k(\mathbf A)\in Q_\varphi.
\]
Space is polynomial in the size of the ordered type structure,
not in the size of \(A\). The bound \(k\) may increase under
renaming and nesting; no sharp bound by the original number of
printed variable names is asserted. Parameters and free relation
variables of inner fixed-point matrices are retained during evaluation.
\end{hblemma}

Under \(\mathsf P=\mathsf{PSPACE}\), ordered capture supplies an
LFP sentence for \(Q_\varphi\), and the pullback lemma removes
the auxiliary ordered structure. This establishes the implication from complexity-class equality in
\cref{thm:p03-w03-p03-059-abiteboul-vianu}; its annex proof also
passes from sentences to queries with their original free variables.
It does not assert that LFP captures all polynomial-time queries on
unordered structures: only the specified PFP query is recovered from
its tuple-type information.


\subsection{Limits of the comparisons}

\begin{hbwarning}{def:p03-w03-p03-060-limits}{Where the inclusions stop}
\begin{enumerate}
\item \emph{Finite structures.} \Cref{thm:p03-w03-p03-060-fp-in-lomega}
      uses that the stages stop at a finite level. On arbitrary structures
      LFP is not contained in \(\mathrm{L}^\omega_{\infty\omega}\): Kolaitis
      reports that well-foundedness is LFP-definable but not
      \(\mathrm{L}_{\infty\omega}\)-definable on all structures, citing
      Dickmann; the inclusion holds again on classes of structures of
      bounded cardinality \cite[p.~54]{Kolaitis2003ExpressivePowerNotes}.
      The Datalog inclusion of \cref{thm:p03-w03-p03-060-datalog-in-existential},
      by contrast, holds on all structures, because the stages of a Datalog
      program stop by stage \(\omega\)
      (\cref{thm:p03-w03-p03-060-datalog-stages}).
\item \emph{Games cannot separate LFP from PFP.} Since both logics lie in
      \(\mathrm{L}^\omega_{\infty\omega}\) on finite structures, no
      \(k\)-pebble argument distinguishes them; Kolaitis makes this point
      before the Abiteboul--Vianu theorem
      \cite[p.~67]{Kolaitis2003ExpressivePowerNotes}.
\item \emph{Proper inclusions.} By
      \cref{thm:p03-w03-p03-060-preservation} and
      \cref{thm:p03-w03-p03-060-datalog-in-existential}, the first-order
      sentence \(\exists x\exists y\,x\neq y\) is not equivalent to any
      Datalog query on finite structures (a homomorphism maps a two-element
      set to a one-element set), and
      \(\forall x\forall y(x\neq y\to E(x,y))\) is not equivalent to any
      Datalog\((\neq)\) query (the inclusion of \(\mathbf K_2\) into
      \(\mathbf K_2\) plus an isolated vertex is a one-to-one
      homomorphism). These are Kolaitis's examples
      \cite[pp.~44 and~46]{Kolaitis2003ExpressivePowerNotes}.
\item \emph{Variable count.} The exact count \(k\) is proved for a single
      fixed point of a first-order formula
      (\cref{thm:p03-w03-p03-060-kv-stages}) and for systems of positive
      \(\exists\mathrm{FO}^k\)-formulas and \(k\)-Datalog\((\neq)\)
      (\cref{thm:p03-w03-p03-060-kv-systems,thm:p03-w03-p03-060-kdatalog}).
      For nested LFP, IFP and PFP formulas only the generous bound of
      \cref{thm:p03-w03-p03-060-fp-in-lomega} is proved.
\item \emph{Infinitary disjunctions are needed.} The least fixed point of a
      system of \(\exists\mathrm{FO}^k\)-formulas is in general not
      first-order (\cref{thm:p03-w03-p03-060-tc-not-fo}), and on the class
      of all finite structures \(k\)-Datalog is properly contained in
      \(\exists\mathrm{L}^\omega_{\infty\omega}\), which defines
      non-computable queries \cite[p.~71]{Kolaitis2003ExpressivePowerNotes}.
\end{enumerate}
\end{hbwarning}

\section{Relationships and sources}
\label{sec:p03-w03-p03-060-relationships}

The existential game is the one-sided form of the \(k\)-pebble game of
\cref{def:p03-w03-p03-056-pebble-game,thm:p03-w03-p03-056-eternal} and of
\cref{def:p03-w03-p03-055-pebble,thm:p03-w03-p03-055-pebble}: a winning
family for the full game is an \(\exists\)-family. The fixed-point
operators and the LFP-into-PFP inclusion are those of
\cref{sec:p03-w03-p03-059-fixed-point-logics}. The zero--one law for
fixed-point logics in \cref{sec:p09-presentation-n-c-002} applies to PFP
through the inclusion of \cref{thm:p03-w03-p03-060-fp-in-lomega}. Datalog
evaluation and its complexity are \cref{def:p09-datalog} and
\cref{sec:p09-presentation-n-c-024}.

\paragraph{Exact sources.} Kolaitis, manuscript of 1 August 2003:
systems of formulas, Definitions~6.12 and~6.15 \cite[pp.~39--40]{Kolaitis2003ExpressivePowerNotes};
Datalog, Datalog\((\neq)\), Examples~6.18, 6.20 and~6.24, Propositions~6.21,
6.23 and~6.26 \cite[pp.~42--46]{Kolaitis2003ExpressivePowerNotes};
partial fixed points and the Gurevich--Shelah remark \cite[pp.~49--51]{Kolaitis2003ExpressivePowerNotes};
Definition~7.1 and Theorem~7.2 \cite[pp.~52--54]{Kolaitis2003ExpressivePowerNotes};
LFP versus PFP \cite[pp.~67--68]{Kolaitis2003ExpressivePowerNotes};
Section~8.1--8.2, Definitions~8.1, 8.4--8.6, Theorems~8.2, 8.3, 8.7 and
Corollary~8.8 \cite[pp.~69--73]{Kolaitis2003ExpressivePowerNotes}.
Gr\"adel, chapter in the web copy: Definition~3.3.7 \cite[p.~155]{Gradel2007FMTDescriptiveChapter};
simultaneous inductions and the Beki\'c principle, Definition~3.3.39,
Lemma~3.3.41 and Theorem~3.3.42 \cite[pp.~171--173]{Gradel2007FMTDescriptiveChapter};
Exercise~3.3.44 and Definition~3.3.45 \cite[p.~174]{Gradel2007FMTDescriptiveChapter};
least versus inflationary fixed points, Theorem~3.3.49 \cite[pp.~175--176]{Gradel2007FMTDescriptiveChapter};
Theorem~3.3.56 \cite[p.~180]{Gradel2007FMTDescriptiveChapter}; Datalog,
Definitions~3.3.57, 3.3.62 and~3.3.63 and Proposition~3.3.64
\cite[pp.~180--182]{Gradel2007FMTDescriptiveChapter}.
